\documentclass[10pt]{article}
\usepackage[T1]{fontenc}
\usepackage[utf8]{inputenc}
\usepackage[french]{babel}
\usepackage{amsmath,amssymb,mathrsfs,mathpazo}
\usepackage{geometry}
\geometry{a4paper,margin=2cm}
\usepackage{graphicx,array,tabularx,tabularray,colortbl,multirow,multicol,booktabs,enumitem,pifont,ulem,fancybox,diagbox,cancel}
\usepackage{pstricks,pstricks-add,pst-plot,pst-tree,pst-node,pst-text,pst-func,pst-math,pst-bspline,pst-3dplot,pst-eucl,pst-solides3d}
\let\pscalculate\undefined
\usepackage{pst-bezier}
\usepackage{tikz,pgfplots,tkz-tab}
\usetikzlibrary{arrows,arrows.meta,patterns,calc,positioning,trees,decorations.markings,shapes.geometric,shapes.misc,shapes.arrows,intersections,decorations.pathmorphing,decorations.pathreplacing,backgrounds,shadows,babel}
\pgfplotsset{compat=1.18}
\usepgfplotslibrary{fillbetween}
\usepackage[np]{numprint}
\usepackage{listings,makecell,pas-tableur,twoopt}
\let\barre\undefined
\usepackage[pstricks]{tablvar}
\definecolor{darkgreen}{RGB}{0,100,0}
\providecommand{\red}{\color{red}}
\providecommand{\blue}{\color{blue}}
\providecommand{\green}{\color{green}}
\providecommand{\black}{\color{black}}
\providecommand{\white}{\color{white}}
\providecommand{\gray}{\color{gray}}
\providecommand{\cyan}{\color{cyan}}
\providecommand{\magenta}{\color{magenta}}
\providecommand{\yellow}{\color{yellow}}
\providecommand{\violet}{\color{violet}}
\providecommand{\orange}{\color{orange}}
\providecommand{\euro}{\texteuro}
\providecommand{\EUR}{\texteuro}
\providecommand{\ieme}{\textsuperscript{e}}
\providecommand{\degres}{\ensuremath{^\circ}}
\providecommand{\degree}{\ensuremath{^\circ}}
\providecommand{\vv}[1]{\overrightarrow{#1}}
\providecommand{\Coord}[1]{#1}
\providecommand{\decosix}{}
\providecommand{\starredbullet}{\ensuremath{\bullet}}
\pagestyle{empty}
\setlength{\parindent}{0pt}
\usepackage{adjustbox}\DeclareUnicodeCharacter{2236}{:}\DeclareUnicodeCharacter{2190}{\ensuremath{\leftarrow}}
\providecommand{\R}{\mathbb{R}}
\providecommand{\N}{\mathbb{N}}
\providecommand{\Z}{\mathbb{Z}}
\providecommand{\Q}{\mathbb{Q}}
\providecommand{\D}{\mathbb{D}}
\providecommand{\C}{\mathbb{C}}
\providecommand{\e}{{\mathrm{e}}}
\providecommand{\vect}[1]{\overrightarrow{#1}}
\providecommand{\vectt}[1]{\overrightarrow{\mathrm{#1}}}
\providecommand{\vvt}[1]{\overrightarrow{\mathrm{#1}}}
\providecommand{\np}[2][]{{#2~#1}}
\providecommand{\vv}[1]{{\overrightarrow{#1}}}
\renewcommand{\d}{{\mathrm{d}}}
\providecommand{\degres}{{^\circ}}
\providecommand{\degre}{{^\circ}}
\providecommand{\blue}{\color{blue}}
\providecommand{\red}{\color{red}}
\providecommand{\enskip}{\quad}
\providecommand{\llbracket}{\lbrack\!\lbrack}
\providecommand{\rrbracket}{\rbrack\!\rbrack}
\providecommand{\euro}{{\text{€}}}
\providecommand{\no}{{\text{n°}}}
\providecommand{\interval}[2]{{[#1;#2]}}
\providecommand{\intervFO}[2]{{[#1;#2[}}
\providecommand{\dotfill}{\ldots}

\providecommand{\ds}{\displaystyle}
\providecommand{\textding}[1]{#1}
\providecommand{\Oij}{\(\mathrm{O};\vec{\imath},\vec{\jmath}\)}
\providecommand{\Oijk}{\(\mathrm{O};\vec{\imath},\vec{\jmath},\vec{k}\)}
\providecommand{\Ouv}{\(\mathrm{O};\vec{u},\vec{v}\)}
\providecommand{\intervFF}[2]{\ensuremath{[#1;#2]}}
\renewcommand{\euro}{\ifmmode\text{\texteuro}\else\texteuro\fi}
\DeclareUnicodeCharacter{00B0}{\ifmmode^{\circ}\else\textdegree\fi}
\DeclareUnicodeCharacter{20AC}{\ifmmode\text{\texteuro}\else\texteuro\fi}
% Énoncés : Éducation nationale. Transcription et corrigés : APMEP, Denis Vergès et contributeurs. Découpage et classement : CoopMaths, Sébastien Lozano, Éric Elter et contributeurs. Adaptation de la présentation : Maths Explorer.
% CC BY-SA : https://coopmaths.fr/www/about/
% https://coopmaths.fr/alea/static/bac/2022/tex/bac_2022_05_sujet2_madagascar_3.tex
\begin{document}
\noindent\textit{APMEP / CoopMaths. CC BY-SA. Présentation adaptée par Maths Explorer.}\par\medskip
\medskip

\emph{Cet exercice est un questionnaire à choix multiples.\\
Pour chacune des questions suivantes, une seule des quatre réponses proposées est exacte.\\
Une réponse fausse, une réponse multiple ou l'absence de réponse à une question ne rapporte ni n'enlève de point.\\
Pour répondre, indiquer sur la copie le numéro de la question et la lettre de la réponse choisie.\\Aucune justification n'est demandée.}

\medskip

\begin{enumerate}
\item Soit $g$ la fonction définie sur $\R$ par $g(x) = x^{\np{1000}} + x$. 

On peut affirmer que :
	\begin{enumerate}
		\item la fonction $g$ est concave sur $\R$.
		\item la fonction $g$ est convexe sur $\R$.
		\item la fonction $g$ possède exactement un point d'inflexion.
		\item la fonction $g$ possède exactement deux points d'inflexion.
	\end{enumerate}
\item On considère une fonction $f$ définie et dérivable sur $\R$. 

\begin{minipage}{0.48\linewidth}
On note $f'$ sa fonction dérivée.

On note $\mathcal{C}$ la courbe représentative de $f$.

On note $\Gamma$ la courbe représentative de $f'$.

On a tracé ci-contre la courbe $\Gamma$.
\end{minipage}\hfill
\begin{minipage}{0.4\linewidth}
\psset{unit=1.2cm}
\begin{pspicture*}(-1.8,-1.6)(3.2,1.4)
\psgrid[gridlabels=0pt,subgriddiv=1,gridwidth=0.15pt]
\psaxes[linewidth=1.25pt,labelFontSize=\scriptstyle]{->}(0,0)(-1.8,-1.6)(3.2,1.4)
\psplot[plotpoints=2000,linewidth=1.25pt,linecolor=red]{-1.8}{3.2}{x 1 add 2.71828 x exp div}
\uput[ur](2,0.45){\red $\Gamma$}
\end{pspicture*}
\end{minipage}

On note $T$ la tangente à la courbe $\mathcal{C}$ au point d'abscisse 0 .

On peut affirmer que la tangente $T$ est parallèle à la droite d'équation : 

\begin{center}
\begin{tabularx}{\linewidth}{*{2}{X}}
\textbf{a.~~}$y =x$ &\textbf{b.~~}$y = 0$\\
\textbf{c.~~}$y = 1$&\textbf{d.~~}$x = 0$
\end{tabularx}
\end{center}
\item On considère la suite $\left(u_n\right)$ définie pour tout entier naturel $n$ par $u_n = \dfrac{(-1)^n}{n+1}$.

On peut affirmer que la suite $\left(u_n\right)$ est : 

\begin{center}
\begin{tabularx}{\linewidth}{*{2}{X}}
\textbf{a.~~}majorée et non minorée.&\textbf{b.~~}minorée et non majorée.\\
\textbf{c.~~}bornée.				&\textbf{d.~~}non majorée et non minorée.
\end{tabularx}
\end{center}

\item Soit $k$ un nombre réel non nul.

Soit $\left(v_n\right)$  une suite définie pour tout entier naturel $n$.

On suppose que $v_0 = k$ et que pour tout $n$, on a $v_n \times v_{n+1} < 0$.

On peut affirmer que $v_{10}$ est :

\begin{center}
\begin{tabularx}{\linewidth}{*{2}{X}}
\textbf{a.~~}positif. &\textbf{a.~~} négatif.\\
\textbf{c.~~}du signe de $k$.&\textbf{d.~~}du signe de $- k$.
\end{tabularx}
\end{center}
\item On considère la suite $\left(w_n\right)$ définie pour tout entier naturel $n$ par : 

\[w_{n+1} = 2w_n - 4\quad \text{et}\quad  w_2 = 8.\]

On peut affirmer que:

\begin{center}
\begin{tabularx}{\linewidth}{*{2}{X}}
\textbf{a.~~}$w_0 = 0$&\textbf{b.~~} $w_0 = 5$.\\
\textbf{c.~~}$w_0 = 10$.&\textbf{d.~~}Il n'est pas possible de calculer $w_0$.
\end{tabularx}
\end{center}

\item \footnote{Cette question ne fait pas partie du sujet donné à Madagascar} On considère la suite $\left(a_n\right)$ définie pour tout entier naturel $n$ par :

\[a_{n+1} = \dfrac{\text{e}^n}{\text{e}^n + 1}a_n\quad \text{et} \quad a_0 = 1.\]

On peut affirmer que :

\begin{center}
\begin{tabularx}{\linewidth}{*{2}{X}}
\textbf{a.~~} la suite $\left(a_n\right)$ est strictement croissante.&\textbf{b.~~}la suite $\left(a_n\right)$ est strictement décroissante.\\
\textbf{c.~~} la suite $\left(a_n\right)$ n'est pas monotone.&\textbf{d.~~} la suite $\left(a_n\right)$ est constante.
\end{tabularx}
\end{center}
\item Une cellule se reproduit en se divisant en deux cellules identiques, qui se divisent à leur tour, et ainsi de suite. 

On appelle temps de génération le temps nécessaire pour qu'une cellule donnée se divise en deux cellules.

On a mis en culture 1 cellule. Au bout de 4 heures, il y a environ \np{4000} cellules.

On peut affirmer que le temps de génération est environ égal à :

\begin{center}
\begin{tabularx}{\linewidth}{*{2}{X}}
\textbf{a.~~}moins d'une minute.&\textbf{b.~~}12 minutes.\\
\textbf{c.~~}20 minutes.		&\textbf{d.~~}1 heure.
\end{tabularx}
\end{center}

\end{enumerate}

\bigskip

\end{document}
